\documentclass[11pt]{article}
\usepackage[margin=1.08in]{geometry}
\usepackage{amsmath,amssymb,amsthm,booktabs}
\usepackage[colorlinks=true,linkcolor=blue,urlcolor=blue,citecolor=blue]{hyperref}
\emergencystretch=3em
\newtheorem{theorem}{Theorem}[section]
\newtheorem{lemma}[theorem]{Lemma}
\newtheorem{proposition}[theorem]{Proposition}
\newtheorem{corollary}[theorem]{Corollary}
\newtheorem{remark}[theorem]{Remark}
\newcommand{\R}{\mathbb R}\newcommand{\C}{\mathbb C}\newcommand{\N}{\mathbb N}
\newcommand{\hist}{ac8979631c1541b059b3c10f590ca174b0b5f6be}
\newcommand{\diagpin}{1bbe6e3ca9869777c205abe88f3303b2dbc07292}
\newcommand{\decl}[2]{\href{https://github.com/lluiseriksson/THE-ERIKSSON-PROGRAMME/blob/\hist/YangMills/OS/#1.lean}{\nolinkurl{#2}}}
\title{Finite Source-Weighted Kernels:\\Strict Peripheral Separation and an Explicit Shrinking-Gap Family}
\author{Lluis Eriksson}
\date{8 November 2026\quad Revision v2}
\begin{document}\maketitle
\begin{abstract}
We specify a finite spin kernel and its source-weight convention, then distinguish
strict separation at each finite size from a bound uniform over sizes and weights.
For $X_L=\{-1,1\}^L$, $\beta\in\R$ and $w:X_L\to(0,\infty)$, the operator has
entries $A(\sigma,\tau)=w(\sigma)\exp(\beta\sum_j\sigma_j\tau_j)$.
It is similar to a real symmetric matrix. A classical positive-matrix argument,
organized through a positive left eigenvector and the pair $|u|-u,|u|+u$,
excludes the negative Perron eigenvalue and gives strict separation of the whole
spectrum. We state the actual specialized historical formal premises rather
than identify a theorem name with an unspecified physical operator.
We also give a fully explicit positive source-weight sequence at each fixed
$\beta>0$ whose relative modulus gap satisfies
$\delta_L\sim2\exp(-2\beta L)$. The construction concentrates weight at the two
constant configurations; its proof uses a rank-two compression and exact
variational bounds. These weights depend globally on the configuration and on
$L$. The result supplies no conclusion for a fixed local-interaction Ising
family. The finite gap, the all-size analytic counterexample and earlier formal
endpoints are distinguished; twelve rational diagnostic cases corroborate the
finite construction but do not prove its limit. Originality of this particular
specialization remains unestablished. No interacting or continuum Yang--Mills
result is claimed.
\end{abstract}

\section{The operator and the question}
Put $X_L=\{-1,1\}^L$, including $L=0$, and let
\[
 b_\beta(s,t)=e^{\beta st},\qquad
 K_{\beta,L}(\sigma,\tau)=\prod_{j=1}^L b_\beta(\sigma_j,\tau_j).
\]
All parameters are finite real numbers. The source-weighted convention here is
\begin{equation}\label{eq:A}
 A_{w,\beta,L}(\sigma,\tau)=w(\sigma)K_{\beta,L}(\sigma,\tau),\qquad
 (Ax)(\sigma)=\sum_{\tau\in X_L}A(\sigma,\tau)x(\tau).
\end{equation}
Thus $A=\operatorname{diag}(w)K$; the weight multiplies rows, not columns.
Every entry is strictly positive for every real $\beta$ when $w>0$.
Writing $D=\operatorname{diag}(\sqrt w)$ and $S=DKD$, we have
\begin{equation}\label{eq:similar}
 A=DSD^{-1},\qquad D^{-1}AD=S.
\end{equation}
The inverses exist because every weight is positive. Since $K$ is real and
symmetric, so is $S$. This is symmetry of a specified finite matrix; entrywise
positivity alone does not assert positive semidefiniteness.

There is no spatial boundary convention hidden in~\eqref{eq:A}: the displayed
product contains temporal sitewise bonds, and any spatial interaction must be
specified separately by $w$. For example an open-chain weight is
$w_{\gamma,L}(\sigma)=\exp(\gamma\sum_{j=1}^{L-1}\sigma_j\sigma_{j+1})$;
it is positive for every real $\gamma$. We do not infer the large-$L$ behavior
of that fixed-$\gamma$ family from the theorem for arbitrary source weights.
At $L=0$ the configuration space has one element, the empty product equals1,
and $A=[w(\emptyset)]$. There is no second eigenvalue at that size.

The generic existence, uniqueness and separation results are classical Perron
matrix theory~\cite{Perron,Ahmadinejad}. The contribution intended here is a
specified interface to the historical formal development, an elementary
organization of its real peripheral proof, and a concrete shrinking-gap family
in the exact kernel class. These are not claims to new Perron theory or to the
first formalization. They are consolidated in this revision of the existing
paper, rather than presented as separate corollary papers.

\section{Finite positive kernels and real peripheral separation}
Let $A$ be a real matrix on a finite nonempty index set, with $A_{ij}>0$.
Classical positive-matrix existence supplies $v>0$ and $Av=\lambda v$.
Taking any coordinate shows $\lambda>0$.

\begin{lemma}[Left vector and subeigenvector equality]\label{lem:left}
There is $\varphi>0$ with $A^{\mathsf T}\varphi=\lambda\varphi$.
If $Au-\lambda u\ge0$ componentwise, then $Au=\lambda u$.
\end{lemma}
\begin{proof}
Apply existence to $A^{\mathsf T}$, obtaining eigenvalue $\nu$.
Pairing its positive eigenvector with $Av$ gives
$(\lambda-\nu)\langle\varphi,v\rangle=0$; the pairing is positive, so
$\nu=\lambda$. For $z=Au-\lambda u\ge0$ the same pairing gives
$\langle\varphi,z\rangle=0$. All positive coefficients $\varphi_i$ force
$z_i=0$. No sign assumption on $u$ is needed for this implication.
\end{proof}

\begin{lemma}[Domination]\label{lem:dom}
Every complex eigenvalue $\mu$ of $A$ obeys $|\mu|\le\lambda$.
\end{lemma}
\begin{proof}
For a nonzero complex eigenvector $x$, let $r=\max_j|x_j|/v_j>0$ and choose
$i$ attaining it. Then $|x_i|=rv_i>0$ and
\[
 |\mu||x_i|=\left|\sum_jA_{ij}x_j\right|
 \le\sum_jA_{ij}|x_j|\le r\sum_jA_{ij}v_j
 =r\lambda v_i=\lambda|x_i|.
\]
Division by $|x_i|$ proves the claim.
\end{proof}

\begin{theorem}[Excluding the negative Perron value]\label{thm:minus}
There is no nonzero real vector $u$ with $Au=-\lambda u$.
Consequently every real eigenvalue $\mu\ne\lambda$ has $|\mu|<\lambda$.
\end{theorem}
\begin{proof}
Suppose $Au=-\lambda u$, and put $h=|u|$ entrywise. The triangle inequality
and $\lambda>0$ give $Ah\ge\lambda h$, so Lemma~\ref{lem:left} gives equality.
Set $p=h-u\ge0$, $q=h+u\ge0$. Then
\[
 Ap=\lambda q,\qquad Aq=\lambda p.
\]
If $p\ne0$, strict positivity of $A$ implies $Ap>0$, and hence $q>0$.
But $q_i=|u_i|+u_i$ vanishes when $u_i\le0$, so $u>0$, which forces
$p=0$, a contradiction. Thus $p=0$ and $h=u\ge0$. The two equations
$Au=\lambda u$ and $Au=-\lambda u$ imply $u=0$. Finally,
Lemma~\ref{lem:dom} and realness leave only $\pm\lambda$ as possible equality
cases, both excluded for a real eigenvalue different from $\lambda$.
\end{proof}

This $p,q$ organization avoids formalizing an equality case of the triangle
inequality. We claim neither priority for the organization nor measured savings
against another formal implementation. It is an interface choice in the
historical development, not a new spectral principle.

\section{The full spectrum of the specified source-weighted operator}
\begin{theorem}\label{thm:whole}
For~\eqref{eq:A}, every complex eigenvalue is real. The Perron eigenvalue is
algebraically simple, and every other eigenvalue has modulus strictly smaller
than it. If $L\ge1$, the finite relative modulus gap
\[
 \delta(A)=1-\max_{\mu\ne\lambda}\frac{|\mu|}{\lambda}
\]
is a well-defined positive number. This assertion supplies no uniform modulus
over unspecified weight sequences.
\end{theorem}
\begin{proof}
Similarity~\eqref{eq:similar} transfers an eigenpair of $A$ to one of the real
symmetric matrix $S$. For $Sy=\mu y$, pairing $y$ with $Sy$ and conjugating
shows $(\mu-\bar\mu)\sum_i|y_i|^2=0$, hence $\mu$ is real.
For a real eigenvalue and a nonzero complex eigenvector of the real matrix $A$,
its real and imaginary parts are real eigenvectors, and at least one is nonzero.
Theorem~\ref{thm:minus} therefore gives the strict modulus assertion for every
non-Perron eigenvalue.

To check simplicity, let $Au=\lambda u$ with $u$ real, and choose
$c=\min_i u_i/v_i$. Then $y=u-cv\ge0$, one coordinate of $y$ is zero, and
$Ay=\lambda y$. If $y\ne0$, strict positivity of $A$ would imply $Ay>0$ at
that zero coordinate, a contradiction. Thus the real Perron eigenspace is
one-dimensional. A real symmetric matrix is diagonalizable, so similarity to
$S$ promotes this geometric assertion to algebraic simplicity for this class.
For $L\ge1$ there are $2^L\ge2$ eigenvalues counted with multiplicity, and the
maximum of the finitely many non-Perron moduli is strictly below $\lambda$.
\end{proof}

The existence of this positive finite gap is quantitative in an existential
sense. The issue left open by qualitative separation is an explicit controlled
modulus as $L$, $w$ and parameters vary. Calling the finite theorem devoid of
all quantitative content would be too strong. Likewise $L=0$ gives a valid
one-dimensional operator and a vacuous separation assertion, not a defined
second-largest eigenvalue. For $\beta=0$, $A=w\mathbf1^{\mathsf T}$ has sole
nonzero eigenvalue $\sum_\sigma w(\sigma)$, so $\delta(A)=1$ for $L\ge1$.
Not every positive-weight family loses its uniform relative gap.

\section{An explicit family with exponentially shrinking relative gap}
\begin{theorem}[Two-corner source weights]\label{thm:corner}
Fix $q\in(0,1)$. For each integer $L\ge1$ put
\[
 H_L(\sigma,\tau)=q^{d_H(\sigma,\tau)},\qquad
 R_L=(1+q)^L,\qquad z_L=q^L,\qquad t_L=\frac{z_L^2}{R_L}.
\]
Assign weight1 to the two constant configurations and weight $t_L^2$ to every
other configuration. Let $D_L=\operatorname{diag}(\sqrt{w_L})$,
$S_L=D_LH_LD_L$, and $A_L=D_L^2H_L$. All weights are strictly positive.
The eigenvalues are real and positive, with a simple largest eigenvalue.
Writing $\delta_L=1-\lambda_2/\lambda_1$ and
$\eta_L=2z_L^2+z_L^4/R_L$, one has
\begin{equation}\label{eq:delta}
 \frac{2z_L-\eta_L}{1+z_L}\ \le\ \delta_L\ \le\
 \frac{2z_L+\eta_L}{1+z_L+\eta_L},\qquad
 \frac{\delta_L}{2q^L}\longrightarrow1.
\end{equation}
For the original kernel~\eqref{eq:A}, take any fixed $\beta>0$ and
$q=e^{-2\beta}$. The same source weights then give
$\delta_L\sim2e^{-2\beta L}$. Thus strict positivity alone cannot imply a
positive volume-uniform relative gap for the arbitrary-source-weight class.
\end{theorem}
\begin{proof}
We omit subscripts temporarily. The Walsh characters
$\chi_J(\sigma)=\prod_{j\in J}\sigma_j$ form an orthogonal basis and satisfy
\[
 H\chi_J=(1-q)^{|J|}(1+q)^{L-|J|}\chi_J.
\]
This identity follows by separating the $L$ sums over $\tau_j$.
Every displayed eigenvalue is positive, so $H$ is positive definite.
Consequently $S=DHD$ is positive definite by congruence, and $A$ is similar to
$S$. All eigenvalues are positive, so the second-largest modulus is exactly
$\lambda_2$. The finite strictness and simplicity follow from
Theorem~\ref{thm:whole}; alternatively they follow from the positive-entry
Rayleigh argument for $S$.

All row and column sums of $H$ equal $R$. Weighted Cauchy--Schwarz gives
\[
 |(Hx)_i|^2\le R\sum_jH_{ij}|x_j|^2,
 \qquad \|Hx\|_2^2\le R^2\|x\|_2^2.
\]
The constant vector attains equality; hence $\|H\|_{2\to2}=R$.
Let $P$ be the orthogonal projection onto the two corner coordinate vectors
and $Q=I-P$. Then $D=P+tQ$. The compression $B=PHP$ has corner block
$\left[\begin{smallmatrix}1&z\\z&1\end{smallmatrix}\right]$ and is zero off
that coordinate subspace. Its two largest eigenvalues are $1+z$ and $1-z$.
The exact decomposition
\[
 S-B=t(PHQ+QHP)+t^2QHQ
\]
therefore yields $\|S-B\|\le(2t+t^2)R=\eta$.
At $L=1$, $P=I$ and $Q=0$; the same statements remain valid.

The normalized sum of the corner vectors has Rayleigh quotient $1+z$ for
$S$, while the perturbation bound gives the upper estimate. Thus
\[
 1+z\le\lambda_1\le1+z+\eta.
\]
Every unit vector in the two-dimensional range of $P$ has Rayleigh quotient
at least $1-z$. That space intersects the orthogonal complement of a top
eigenvector nontrivially, proving $\lambda_2\ge1-z$.
On the orthogonal complement of the corner-sum vector, $B$ has Rayleigh
quotient at most $1-z$, and $S$ has quotient at most $1-z+\eta$.
The finite-dimensional min-max principle consequently gives
\[
 1-z\le\lambda_2\le1-z+\eta.
\]
Dividing by the appropriate positive bounds on $\lambda_1$ proves both
inequalities in~\eqref{eq:delta}. Finally
$\eta_L/z_L=2z_L+z_L^3/R_L\to0$, so both bounds divided by $2z_L$ tend to1.
Since $K_{\beta,L}=e^{\beta L}H_L$ when $q=e^{-2\beta}$, the positive scalar
factor changes no eigenvalue ratio and proves the claim for~\eqref{eq:A}.
\end{proof}

\paragraph{Scope of this counterexample.}
The source weights depend on $L$ and globally distinguish the two constant
configurations. Their minimum tends to zero. They are not asserted to come from
one fixed local nearest-neighbour coupling $\gamma$. The theorem is an actual
family in the source-weighted kernel class, not merely a changing $2\times2$
example, but it does not settle the fixed-$(\beta,\gamma)$ Ising case or exclude
uniform estimates under stronger separately proved weight restrictions.
The counterexample removes no hypothesis from the physical Yang--Mills problem.

\section{Historical formal endpoint and finite diagnostics}
The historical formal source is the immutable checkpoint
\begin{center}\small\texttt{\hist}\end{center}
It declares \texttt{leanprover/lean4:v4.29.0-rc6}; its Mathlib revision is
\begin{center}\small\texttt{07642720480157414db592fa85b626dafb71355b}.\end{center}
No Lean/Lake build or axiom oracle was replayed for this revision.
The broad build and oracle totals printed in the earlier artifact are not
reasserted as newly authenticated evidence.
The formal configuration type \texttt{Fin L} $\to$ \texttt{Fin 2} is identified
with $X_L$ by mapping the site values $0\mapsto+1$, $1\mapsto-1$.
The definitions are \decl{Z2Identification}{z2Bond},
\decl{SpatialExtent}{spatialKernel}, and
\decl{PerronKernel}{sourceWeightedKernelL}. The positive source-weight premise
is explicit; \decl{PerronKernel}{vacuum_exists_of_sourceWeight} supplies the
positive Perron vector used by the separation endpoint.

For clarity the full mathematical signature of
\decl{PerronGap}{coupled_gap_all_eigenvalues} is as follows.
Take $L\in\N$, $w:X_L\to\R$, $\beta\in\R$, $v:X_L\to\R$,
$\lambda\in\R$, $x:X_L\to\C$, $\sigma_1\in X_L$, and $\mu\in\C$.
Its premises are
\begin{align*}
 &\forall\sigma,\quad w(\sigma)>0,\qquad v(\sigma)>0,\\
 &\forall\sigma,\quad\sum_\tau A_{w,\beta,L}(\sigma,\tau)v(\tau)
       =\lambda v(\sigma),\\
 &x(\sigma_1)\ne0,\\
 &\forall\sigma,\quad\sum_\tau
       (A_{w,\beta,L}(\sigma,\tau):\C)x(\tau)=\mu x(\sigma),\\
 &\mu\ne(\lambda:\C).
\end{align*}
Its conclusion is $\|\mu\|<\lambda$.
The declaration \decl{PerronGap}{sourceWeighted_eq_sym} gives precisely
$A=DSD^{-1}$, and \decl{PerronGap}{coupled_eigenvalue_real} gives realness.
The final endpoint explicitly takes real or imaginary parts to recover a
nonzero real eigenvector. This is source-level correspondence, not new formal
verification of the endpoint, of Theorem~\ref{thm:corner}, or of every analytic
claim in this paper. In particular the algebraic-simplicity argument above is
not represented as a newly checked Lean declaration.

\paragraph{Finite exact checks of the new construction.}
At diagnostic source checkpoint
\begin{center}\small\texttt{\diagpin}\end{center}
\nolinkurl{scripts/perron_two_corner_diagnostic_210.py} checks $q=1/3,1/2,2/3$
and $L=1,2,3,4$ using rational arithmetic. It checks the source-weight
similarity, row/column sums, the rank-two decomposition and exact rational
LDL certificates for positive definiteness and the two upper variational
bounds. There are twelve distinct parameter cases, 1,020 entry checks and
48 matrix certificates; normal and optimized execution agree. Four negative
checks reject the excluded endpoints $q=0,1$, an incorrect top-eigenvalue bound
and the wrong column-weight convention. These are finite diagnostic checks,
not a proof of the all-$L$ asymptotic, which is proved analytically above.
The full results, original inputs, actual child exit codes and archive hashes
are retained in the revision's reproducibility package. Execution used an
owned Colab Linux CPU runtime, which was stopped after the run; billed units
were not available. No floating numerical spectrum is used as evidence.

\section{Correction record, originality and review context}
The earlier exact PDF did not specify its full source-weight convention and
specialized formal premises. Sections1 and5 now supply them. The earlier
numerical table is removed: its cited19-line open-chain transcript at the
stated checkpoint does not match the printed ratios and contains no $L=5$
row. For example, at $\beta=0.8$, $\gamma=1.2$, $L=2$ the old table prints
$0.9205$, whereas the cited transcript prints $0.908247$. This establishes a
provenance discrepancy, not impossibility of those numbers for some different,
unverified convention. Both the old artifact and the full discrepancy record
are preserved. Neither a historical finite grid nor a comparison with an Ising
critical line is used to infer a limit here.

The ingredients of the matrix arguments are classical. The exact source-weight
family in Theorem~\ref{thm:corner} is an explicit specialization of elementary
compression and perturbation reasoning; a targeted primary-source comparison
did not establish its priority. Absence of a matching search result is not
novelty evidence. Perron's original publication details are verified at its
publisher; its full original proof was not obtained. The primary research
source~\cite{Ahmadinejad}, Section2.3, explicitly treats the established Perron
framework; no algorithmic or field-level novelty from that work is attributed
to this paper. Further relevance of the formal interface requires practical
reuse, and no comparative formalization effort has been measured.

\paragraph{Model assessments and review context.}
Codex assisted this manuscript's preparation, source audit and exact diagnostic
orchestration. The preparing workflow is involved in the artifact; independent
review or verified context isolation is not asserted. Exact served backend and
reasoning effort were unexposed. Reports for the historical PDF retain their
own referent. Any numerical assessment of this revision must be recorded
separately against its exact frozen PDF hash; no score from another artifact
is applied. AI assessment, historical source inspection, finite exact checking,
human peer review, Working-paper deposit and editorial acceptance are distinct.
No human acceptance or model admission quorum is asserted here.

\begin{thebibliography}{9}
\bibitem{Perron}O.~Perron, \emph{Zur Theorie der Matrices},
Mathematische Annalen \textbf{64} (1907), 248--263.
\href{https://doi.org/10.1007/BF01449896}{doi:10.1007/BF01449896}.
\bibitem{Ahmadinejad}A.~Ahmadinejad, A.~Jambulapati, A.~Saberi and A.~Sidford,
\emph{Perron-Frobenius Theory in Nearly Linear Time: Positive Eigenvectors,
M-matrices, Graph Kernels, and Other Applications}, arXiv:1810.02348v1 (2018),
Section2.3. \url{https://arxiv.org/abs/1810.02348}.
\bibitem{Lean}L.~de~Moura and S.~Ullrich,
\emph{The Lean4 theorem prover and programming language}, CADE-28 (2021),625--635.
\bibitem{Mathlib}The mathlib Community, \emph{The Lean mathematical library},
CPP 2020, 367--381. Exact historical revision is pinned in Section5.
\end{thebibliography}
\end{document}

